\newif\ifuseseminar
\useseminartrue
\input{../../latex_common/header.tex.frag}

\title{Prova de Física Matemática II -- Distribuições, Funções de Green e EDO (Recuperação)}

\begin{document}

\begin{enumerate}

	\item Considere a função degrau de Heaviside $\theta(x)$ e a distribuição
	      delta de Dirac $\delta(x)$.
	      \begin{enumerate}
		      \item Partindo do Teorema Fundamental do Cálculo, mostre que
		            $$
			            \frac{\dd}{\dd x}\theta(x-x_0) = \delta(x-x_0).
		            $$
		      \item Usando o resultado do item anterior, calcule
		            $\dfrac{\dd}{\dd x}\big[e^{-x}\theta(x)\big]$ e
		            $\dfrac{\dd^2}{\dd x^2}\big[e^{-x}\theta(x)\big]$, usando a
		            propriedade $g(x)\delta(x-a) = g(a)\delta(x-a)$. Interprete
		            fisicamente o aparecimento do termo em $\delta'(x)$ na
		            segunda derivada.
		      \item Considere a lorentziana
		            $$
			            L_\eta(x) = \frac{1}{\pi}\,\frac{\eta}{x^2+\eta^2}.
		            $$
		            Mostre que $L_\eta(x)$ é normalizada para qualquer $\eta>0$
		            e que
		            $$
			            \lim_{\eta\to0}\int_{-\infty}^{\infty} L_\eta(x)\, f(x)\,\dd x = f(0).
		            $$
	      \end{enumerate}

	\item Considere o operador $\dfrac{\dd^2}{\dd x^2}-k^2$ ($k>0$) atuando
	      sobre funções definidas em toda a reta real.
	      \begin{enumerate}
		      \item Verifique, por substituição direta, que
		            $$
			            G(x,x_0) = -\frac{1}{2k}\, e^{-k|x-x_0|}
		            $$
		            satisfaz $\left(\dfrac{\dd^2}{\dd x^2}-k^2\right)G(x,x_0) = \delta(x-x_0)$.
		      \item Mostre que $G(x,x_0)$ não é única: somando uma solução
		            homogênea $c_1 e^{kx}+c_2 e^{-kx}$ à solução, a equação
		            continua satisfeita. Explique por que, ao contrário do caso
		            $k=0$, exigir que $G$ seja limitada quando $x\to\pm\infty$
		            já basta para fixar $c_1=c_2=0$.
	      \end{enumerate}

	\item Considere a equação de Bessel de ordem $3$:
	      $$
		      x^2y'' + xy' + (x^2-9)y = 0.
	      $$
	      \begin{enumerate}
		      \item Classifique o ponto $x_0=0$, determine o polinômio
		            indicial e suas raízes.
		      \item Encontre a solução correspondente ao maior índice da
		            equação indicial usando o método de Frobenius (relação de
		            recorrência).
		      \item Classifique o tipo da solução linearmente independente
		            (Caso 1, Caso 2-A, Caso 2-B-i ou Caso 2-B-ii) e discuta o
		            aparecimento de termos logarítmicos.
	      \end{enumerate}

	\item Considere a equação diferencial
	      $$
		      xy'' + (3-x)y' + \lambda y = 0.
	      $$
	      \begin{enumerate}
		      \item Resolva a equação em torno de $x_0=0$ por séries de
		            potências e obtenha a relação de recorrência.
		      \item Determine para quais valores de $\lambda$ a solução se
		            reduz a um polinômio.
	      \end{enumerate}

\end{enumerate}

\end{document}
